\section{Topological phase kernel and generation of the seed catalog}
\label{tpk:capitulo}

The two APP evaluations acquire generative efficacy when it is specified
which evaluation acts, at which position it is performed, how that position moves
and what information each transition preserves. The \emph{Topological Phase Kernel},
abbreviated TPK and called in Spanish \emph{núcleo topológico de fase}, brings
these operations together in dynamics on enriched states. Its definition
comprises selection, transport, arithmetic updating and preservation of the
record. The finite construction developed in this chapter provides
the seeds on which the regional lifts and the coinductive
extension will act.

The main result of this stage is a generated and classified catalog:
the initial conditions of the two cursors produce integer emissions;
their ternary reading determines oriented words; the dihedral action organizes
these words into orbits. Each step preserves its preimages and the relations
that allow a return to the preceding description. The chain of cardinalities
\[
 104\,976\longrightarrow468\longrightarrow243\subset729,
 \qquad 243\longrightarrow43,
\]
abbreviates different objects. Each object is constructed below and each
count is justified, before any of them is used as data for
regional selection.

\subsection{Phase selection, transport and enriched state}
\label{tpk:operaciones}

\subsubsection{Tritic selector of operational phase}

The nonadic phase takes its values in \(D_9=\{1,\ldots,9\}\). The tritic
selector is the function
\begin{equation}
 M_{\mathrm{ph}}:D_9\longrightarrow\{-1,0,+1\},\qquad
 M_{\mathrm{ph}}(g)=
 \begin{cases}
 +1,&g\in\{1,4,7\},\\
 -1,&g\in\{2,5,8\},\\
 0,&g\in\{3,6,9\}.
 \end{cases}
 \label{tpk:selector}
\end{equation}
In the emitting realization, the value \(+1\) activates the additive evaluation,
\(-1\) activates the multiplicative evaluation and \(0\) records the neutral
event. The cardinal orientation of the cursor is another coordinate of the state.
This separation allows both the operation being read and the
direction in which its evaluation point is transported to be fixed simultaneously.

For transitions numbered by \(t\geq1\), the phase preceding the reading is
\begin{equation}
 g_t=1+[(t-1)]_9,\qquad \epsilon_t=M_{\mathrm{ph}}(g_t),
 \label{tpk:fase}
\end{equation}
where \([n]_b\) denotes the representative of \(n\) modulo \(b\) in
\(\{0,\ldots,b-1\}\). A window of nine transitions contains the sequence
\((+1,-1,0,+1,-1,0,+1,-1,0)\).

\subsubsection{Oriented cursors and cardinal transport}

Indices \(I_9=\{0,\ldots,8\}\) are used for positions and positive
labels \((a,b)=(i+1,j+1)\) for APP evaluations. A finite cursor is
\[
 c=(i,j;d)\in\mathcal S:=I_9^2\times\mathcal D,
 \qquad \mathcal D=\{N,E,S,O\}.
\]
The four directions act through
\[
 v_N=(-1,0),\quad v_E=(0,1),\quad
 v_S=(1,0),\quad v_O=(0,-1).
\]
The cardinal map corresponding to \(d\) is
\[
 T_d(i,j)=\bigl([i+(v_d)_1]_9,[j+(v_d)_2]_9\bigr).
\]
Thus, \(M_{\mathrm{ph}}\) and \(T_d\) have different domains, codomains and
actions. Tritic activation selects one of the two cursors;
cardinal transport modifies its position in the direction that cursor
already preserves.

The integer lift of the cursor is
\(\widetilde c=(x,y;d)\in\mathbb Z^2\times\mathcal D\), with finite position
\(([x]_9,[y]_9)\). Its turn counts are recovered in
\begin{equation}
 x=9\lfloor x/9\rfloor+[x]_9,\qquad
 y=9\lfloor y/9\rfloor+[y]_9.
 \label{tpk:vueltas}
\end{equation}
These two divisions concern spatial transport. The division of an
additive or multiplicative evaluation constitutes a distinct arithmetic
register, which will be preserved alongside them.

\subsubsection{Composition of the transition}

The state of the finite realization includes the two lifted cursors,
their directions, the phase, the accumulators of the active window, the ordered
list of blocks already emitted and the transition register. Each local
register contains the preceding positions, the integer evaluations of both
sheets, their residues and quotients, the active regime and the orientations.
The prefix and boundary fields that appear in a prolongation will be
incorporated into the same state with their compatibility rules.

The transition is ordered as
\begin{equation}
 \mathcal U_t=\mathsf{Mem}_t\circ\mathsf{Tra}_t\circ\mathsf{Sel}_t.
 \label{tpk:composicion}
\end{equation}
Selection determines the regime and preserves the sample prior to any
displacement. Transport applies the active movement and the cardinal inversion
prescribed by the calendar. Updating incorporates that sample
into the accumulators and the register; upon completing a window, it emits its block.
The composition therefore uses a sample taken before transport even though
its definitive archiving occurs at the end of the transition.

\subsection{Initial conditions and reversible enumeration}
\label{tpk:semillas}

The additive and multiplicative sheets have independent cursors. The
domain of initial conditions is the ordered product
\[
 \Omega_{\mathrm{seed}}=\mathcal S_+\times\mathcal S_\times,
 \qquad |\mathcal S|=9^2\cdot4=324.
\]
Consequently,
\begin{equation}
 |\Omega_{\mathrm{seed}}|=324^2=104\,976.
 \label{tpk:censo-inicial}
\end{equation}
The number \(81\) counts positions; \(324\) counts oriented positions
of one cursor; \(104\,976\) counts ordered pairs of cursors. Choosing
a position in APP preserves only part of an initial condition
of the emitter.

The numbering is fixed
\[
 \operatorname{ind}(N)=0,\quad\operatorname{ind}(E)=1,\quad
 \operatorname{ind}(S)=2,\quad\operatorname{ind}(O)=3,
\]
and one defines
\begin{equation}
 \nu(i,j;d)=4(9i+j)+\operatorname{ind}(d).
 \label{tpk:indice-semilla}
\end{equation}

\begin{proposicion}[Exact enumeration of initial conditions]
\label{tpk:enumeracion}
The map \(\nu\) is a bijection from \(\mathcal S\) onto
\(\{0,\ldots,323\}\). The index pair is encoded reversibly
by \(324\nu_++\nu_\times\in\{0,\ldots,104\,975\}\).
\end{proposicion}
\begin{proof}
Division by four recovers
\[
 \operatorname{ind}(d)=[\nu]_4,\qquad
 j=[\lfloor\nu/4\rfloor]_9,\qquad i=\lfloor\nu/36\rfloor.
\]
The domain bounds guarantee that these three quantities belong to
their index sets and reconstruct \(\nu\) through
\eqref{tpk:indice-semilla}. For the pair, Euclidean division by
\(324\) recovers quotient \(\nu_+\) and residue \(\nu_\times\).
\end{proof}

The enumeration traverses the complete domain before any regional
grouping. The indices are used to locate the preimages of an
emission, keeping explicit the condition that produced each trajectory.

\subsection{Evaluation of the sheets and arithmetic register}
\label{tpk:lecturas}

At a positive position \((a,b)\), the integer evaluations are
\(A_+(a,b)=a+b\) and \(A_\times(a,b)=ab\). For \(n>0\), one maintains
\begin{equation}
 \rho_9(n)=1+[n-1]_9,\qquad
 q_9(n)=\lfloor(n-1)/9\rfloor,\qquad
 n=\rho_9(n)+9q_9(n).
 \label{tpk:division-evaluacion}
\end{equation}
Additive activation incorporates \(\rho_9(A_+)\) into the corresponding
accumulator. Multiplicative activation uses an observable of the
residue \(\rho_9(A_\times)\), defined below.

\subsubsection{Discrete exponent in the units modulo nine}

The powers of \(2\) modulo \(9\) traverse
\[
 2^0,2^1,\ldots,2^5\equiv1,2,4,8,7,5\pmod9,
 \qquad 2^6\equiv1\pmod9.
\]
The discrete logarithm in this group of six units is represented by
\begin{equation}
 \ell_9(1)=0,\quad\ell_9(2)=1,\quad\ell_9(4)=2,
 \quad\ell_9(8)=3,\quad\ell_9(7)=4,\quad\ell_9(5)=5.
 \label{tpk:logaritmo}
\end{equation}
The emission law extends the observable through
\(\ell_9(3)=\ell_9(6)=\ell_9(9)=0\). To preserve its provenance, the
recorded sample includes the original residue and the indicator of membership
in the group of units. In this way, an observed value of zero can come
from the unit \(1\) or from one of the three noninvertible residues, and its
origin remains recoverable.

\subsubsection{Order of reading and movement}

Evaluations are performed at the positions preceding the transition.
If \(\epsilon_t=+1\), the additive reading is accumulated and its cursor is moved.
If \(\epsilon_t=-1\), the multiplicative discrete exponent is accumulated and the
second cursor is moved. If \(\epsilon_t=0\), the neutral counter is
incremented and both positions are preserved.

At the end of steps \(27\) and \(54\), both directions are replaced
by their opposites. The windows are chained together while preserving positions and
orientations. At the start of each window only its three local
accumulators are reset to zero; the previous register remains available.

\begin{proposicion}[Inversion of lifted transport]
\label{tpk:inversion}
Let \(a,b\in\{0,1\}\) be the activation and inversion indicators of a
cursor. If \(\operatorname{opp}\) interchanges \(N\leftrightarrow S\) and
\(E\leftrightarrow O\), the transformation
\[
 F_{a,b}(z,d)=\bigl(z+a v_d,\operatorname{opp}^{\,b}(d)\bigr)
\]
has an inverse
\[
 F_{a,b}^{-1}(z',d')=
 \bigl(z'-a v_{\operatorname{opp}^{\,b}(d')},
       \operatorname{opp}^{\,b}(d')\bigr).
\]
Spatial reduction modulo nine intertwines this transformation with
finite transport.
\end{proposicion}
\begin{proof}
Applying opposition twice recovers the initial direction. Once
that direction is recovered, subtracting \(a v_d\) undoes the displacement.
For spatial projection, the equality
\([z+a v_d]_9=[[z]_9+a v_d]_9\) suffices, component by component. Direction inversion
preserves this equality.
\end{proof}

The phase determines the indicators used by the inverse. The spatial
and arithmetic registers can thus be preserved together, even when
crossing the boundary of the representative \(I_9^2\).

\subsection{Six windows and decimal emission by blocks}
\label{tpk:emision}

\subsubsection{Accumulation in a nonadic window}

Window \(m\), with \(1\leq m\leq6\), comprises the steps
\(9m-8,\ldots,9m\). Let \(S_m\) be the sum of its active positive additive
residues, \(E_m\) the sum of its active discrete exponents and
\(Z_m\) its number of neutral events.

\begin{lema}[Activation multiplicities]
\label{tpk:tres-eventos}
Each window contains three additive activations, three multiplicative activations and
three neutral activations. In particular, \(Z_m=3\).
\end{lema}
\begin{proof}
The phases of a window traverse exactly \(1,\ldots,9\). Each
of the three sets of \eqref{tpk:selector} has cardinality three.
\end{proof}

The six windows form \(54\) transitions and produce six blocks.
This length belongs to the initial emitter defined here. Continuation
in depth subsequently acts on the words and their compatible registers.

\subsubsection{Decimal encoding rule}

The emission law specifies the digits
\begin{align}
 d_{m,1}&=[S_m]_{10},\notag\\
 d_{m,2}&=[S_m+E_m]_{10},\notag\\
 d_{m,3}&=[3S_m+5E_m+7Z_m]_{10},\notag\\
 U_m&=100d_{m,1}+10d_{m,2}+d_{m,3}.
 \label{tpk:bloque-decimal}
\end{align}
The weights \(100,10,1\) are the hundreds, tens and units positions.
The subscript \(10\) indicates the decimal residue, and each block satisfies
\(0\leq U_m\leq999\). The emitted word is the ordered sequence
\(\mathbf U=(U_1,\ldots,U_6)\), with a fixed width of three digits per block.

The coefficients \(3,5,7\) constitute the explicit weights of this emission
law. Once the rule is fixed, each digit is calculated from the
accumulators, and the results of this chapter follow from that rule.
The number \(1\) appearing in its reduced form has a specific additional
provenance: \(7Z_m=21\equiv1\pmod{10}\).

If \(s_m=[S_m]_{10}\) and \(e_m=[E_m]_{10}\), one obtains
\begin{equation}
 G(s_m,e_m)=100s_m+10[s_m+e_m]_{10}
                   +[3s_m+5e_m+1]_{10}=U_m.
 \label{tpk:acoplamiento-firmas}
\end{equation}
The letter \(e_m\) here denotes a component of the discrete-exponent
signature. Its definition belongs to the finite emitter.

\begin{proposicion}[Recovery of the two signatures]
\label{tpk:recuperacion-firmas}
The map \(G:\{0,\ldots,9\}^2\to\{0,\ldots,999\}\) is injective.
For a block \(U=100a+10b+c\) in its image, the inverse is
\[
 s=a,\qquad e=[b-a]_{10},\qquad
 c=[3s+5e+1]_{10}.
\]
Consequently, the word \(\mathbf U\) determines the two six-component
signatures \(\mathbf s\) and \(\mathbf e\).
\end{proposicion}
\begin{proof}
The last two summands of \eqref{tpk:acoplamiento-firmas} belong to
\(\{0,\ldots,99\}\), so the hundreds digit recovers \(s\). The tens digit
is \([s+e]_{10}\); subtracting \(s\) and reducing modulo ten recovers
\(e\). The units digit verifies the third congruence. Reconstruction is
applied independently to the six positions of the word.
\end{proof}

The integer totals \(S_m,E_m\) are recovered from the recorded readings;
the hundreds and tens digits recover their residues. For example, a total
\(S_m=15\) produces \(s_m=5\). This difference identifies exactly
what information each level of encoding preserves.

\subsubsection{Positional encoding of the six-block word}

The six blocks admit a fixed-width integer encoding,
\begin{equation}
 \mathscr N(\mathbf U)=\sum_{m=1}^{6}U_m1000^{6-m},
 \qquad0\leq\mathscr N(\mathbf U)<1000^6.
 \label{tpk:entero-base1000}
\end{equation}
The first block occupies the highest-weight position and the sixth, the units
position in base one thousand. A block \(058\) occupies a complete position with integer
value \(58\); its three-digit width is part of the decimal
presentation of that position.

\begin{proposicion}[Positional recovery of emissions]
\label{tpk:inversa-base1000}
The encoding \eqref{tpk:entero-base1000} determines the complete word
through
\[
 U_m=\left[\left\lfloor
 \frac{\mathscr N(\mathbf U)}{1000^{6-m}}\right\rfloor\right]_{1000},
 \qquad1\leq m\leq6.
\]
\end{proposicion}
\begin{proof}
Upon division by \(1000^{6-m}\), the earlier blocks produce integer
multiples of \(1000\), block \(m\) produces \(U_m\), and the contribution
of later blocks belongs to \([0,1)\). The integer part eliminates this
last contribution and the residue modulo \(1000\) eliminates the earlier ones.
\end{proof}

This encoding preserves the six emissions. Its domain has ordered
positions in base one thousand, whereas the APP domain has spatial positions
modulo nine. The composition maintains both indices: \(m\) identifies an
emission window and \((i,j)\) identifies the cell visited during a
transition in that window.

\subsection{Complete calculation of an emission beginning with 501}
\label{tpk:ejemplo501}

The initial conditions are fixed
\[
 c_{+,0}=(0,4;N),\qquad c_{\times,0}=(2,3;N),
 \qquad \nu_+=16,\quad\nu_\times=84.
\]
The initial positive positions are \((1,5)\) and \((3,4)\). The preceding
chronological index is \(t=0\), the first phase read is \(g_1=1\) and the
register is empty. These coordinates are the inputs
to the calculation; the blocks are obtained by executing the recurrence.

\begin{tablacompleta}[htbp]
\caption{First window: preceding readings and active accumulation.}
\label{tpk:tabla-501}
\small
\begin{tabular}{@{}rcccrrr@{}}
\toprule
\(t\)&\(\epsilon_t\)&position \(+\)&position \(\times\)&
\(\Delta S\)&\(\Delta E\)&\((S,E,Z)\)\\
\midrule
1&\(+1\)&\((1,5)\)&\((3,4)\)&6&0&\((6,0,0)\)\\
2&\(-1\)&\((9,5)\)&\((3,4)\)&0&0&\((6,0,0)\)\\
3&0&\((9,5)\)&\((2,4)\)&0&0&\((6,0,1)\)\\
4&\(+1\)&\((9,5)\)&\((2,4)\)&5&0&\((11,0,1)\)\\
5&\(-1\)&\((8,5)\)&\((2,4)\)&0&3&\((11,3,1)\)\\
6&0&\((8,5)\)&\((1,4)\)&0&0&\((11,3,2)\)\\
7&\(+1\)&\((8,5)\)&\((1,4)\)&4&0&\((15,3,2)\)\\
8&\(-1\)&\((7,5)\)&\((1,4)\)&0&2&\((15,5,2)\)\\
9&0&\((7,5)\)&\((9,4)\)&0&0&\((15,5,3)\)\\
\bottomrule
\end{tabular}
\notatabla{The positions are positive labels of the cells before
movement. Neutral events increment only \(Z\).}
\end{tablacompleta}

The three active additive evaluations are
\[
 1+5=6,\qquad9+5=14=5+9,\qquad8+5=13=4+9,
\]
so that \(S_1=6+5+4=15\). The three active multiplicative evaluations
are \(3\cdot4=12\), \(2\cdot4=8\) and \(1\cdot4=4\). Their positive
residues are \(3,8,4\), with observables \(0,3,2\); therefore,
\(E_1=5\). Finally \(Z_1=3\). The decimal rule gives
\[
 d_{1,1}=[15]_{10}=5,\quad
 d_{1,2}=[15+5]_{10}=0,\quad
 d_{1,3}=[45+25+21]_{10}=1,
\]
and one obtains \(U_1=501\).

The second additive reading above has arithmetic quotient \(1\),
while its lifted cursor is at \((-1,4;N)\), with vertical turn
count \(-1\). These are two different data in the same register. Both
enter into the full recovery of the evaluation and its position.

\begin{tablacompleta}[htbp]
\caption{The six windows of the initial condition \((16,84)\).}
\label{tpk:tabla-seis501}
\begin{tabular}{@{}rrrrrrrr@{}}
\toprule
\(m\)&steps&\(S_m\)&\(E_m\)&\(Z_m\)&\(s_m\)&\(U_m\)&\([-U_m]_3\)\\
\midrule
1&1--9&15&5&3&5&501&0\\
2&10--18&6&5&3&6&614&1\\
3&19--27&24&5&3&4&498&0\\
4&28--36&21&5&3&1&169&2\\
5&37--45&12&5&3&2&272&1\\
6&46--54&12&5&3&2&272&1\\
\bottomrule
\end{tabular}
\end{tablacompleta}

The resulting word and its signatures are
\begin{align*}
 \mathbf U&=(501,614,498,169,272,272),\\
 \mathbf s&=(5,6,4,1,2,2),&
 \mathbf e&=(5,5,5,5,5,5).
\end{align*}
The cardinal inversion after step \(27\) modifies the paths of the
last three windows. The repeated block \(272\) corresponds to two
successive windows with their own positions and chronological register.

As a second case, the minimal initial condition \((\nu_+,\nu_\times)=(0,0)\)
produces
\[
 \mathbf U=(252,109,252,903,850,850),\quad
 \mathbf s=(2,1,2,9,8,8),\quad
 \mathbf e=(3,9,3,1,7,7).
\]
In this case, both cursors return to their initial position and orientation
after \(54\) transitions. The nonadic phase completes six turns, the
chronological counter advances by \(54\) units and the register contains
\(54\) samples. The positional reading of the return and
the recorded history describe distinct properties of the execution.

\subsection{Factorization of the finite census}
\label{tpk:censo}

The signatures \(f_+(\nu)=\mathbf s\) and \(f_\times(\nu)=\mathbf e\) are
calculated by executing the corresponding cursor during the six windows.
The independence of positions and directions between the two cursors gives
\begin{equation}
 T_{\mathrm{em}}(\nu_+,\nu_\times)
   =G^{(6)}\bigl(f_+(\nu_+),f_\times(\nu_\times)\bigr),
 \label{tpk:factorizacion}
\end{equation}
where \(G^{(6)}\) applies \eqref{tpk:acoplamiento-firmas} at each position.
The factorization allows \(324+324\) trajectories to be enumerated first and
the distinct signatures to be combined afterward, preserving the multiplicity of
their preimages.

\subsubsection{Images of the two signatures}

The following tables include all signatures. Their multiplicity is the
number of initial conditions producing them; the minimum index locates
one of those conditions. The complete preimage is defined by
\[
 f_\sigma^{-1}(a)=\{\nu\in\{0,\ldots,323\}:f_\sigma(\nu)=a\}.
\]
The recurrence of the preceding sections determines each element of this
set through a finite integer calculation.

\begin{tablacompleta}[htbp]
\caption{The eighteen additive signatures.}
\label{tpk:firmas-aditivas}
\begin{tabular}{@{}lrr@{\qquad}lrr@{}}
\toprule
signature&multiplicity&\(\nu_{\min}\)&signature&multiplicity&\(\nu_{\min}\)\\
\midrule
\texttt{122564}&18&17&\texttt{564122}&18&16\\
\texttt{122898}&18&24&\texttt{645212}&18&4\\
\texttt{212645}&18&5&\texttt{654889}&18&33\\
\texttt{212988}&18&0&\texttt{889221}&18&13\\
\texttt{221456}&18&29&\texttt{889654}&18&32\\
\texttt{221889}&18&12&\texttt{898122}&18&25\\
\texttt{456221}&18&28&\texttt{898465}&18&20\\
\texttt{465898}&18&21&\texttt{988212}&18&1\\
\texttt{546988}&18&9&\texttt{988546}&18&8\\
\bottomrule
\end{tabular}
\end{tablacompleta}

\begin{tablacompleta}[htbp]
\caption{The twenty-six multiplicative signatures.}
\label{tpk:firmas-multiplicativas}
\begin{tabular}{@{}lrr@{\qquad}lrr@{}}
\toprule
signature&multiplicity&\(\nu_{\min}\)&signature&multiplicity&\(\nu_{\min}\)\\
\midrule
\texttt{000000}&108&8&\texttt{555555}&24&84\\
\texttt{177393}&8&1&\texttt{555717}&8&14\\
\texttt{177555}&8&54&\texttt{555771}&8&18\\
\texttt{177771}&8&11&\texttt{555933}&8&24\\
\texttt{339555}&8&6&\texttt{717555}&8&12\\
\texttt{339771}&8&15&\texttt{717717}&8&21\\
\texttt{339933}&8&59&\texttt{717933}&8&19\\
\texttt{393177}&8&0&\texttt{771177}&8&9\\
\texttt{393393}&8&45&\texttt{771339}&8&13\\
\texttt{393555}&8&40&\texttt{771555}&8&16\\
\texttt{555177}&8&52&\texttt{933339}&8&57\\
\texttt{555339}&8&4&\texttt{933555}&8&26\\
\texttt{555393}&8&41&\texttt{933717}&8&17\\
\bottomrule
\end{tabular}
\end{tablacompleta}

\begin{proposicion}[Image and multiplicities of the emitter]
\label{tpk:468}
The image \(\mathcal U_6=\operatorname{im}T_{\mathrm{em}}\) contains
\(468\) distinct words. Their preimages are distributed among \(432\)
fibers of cardinality \(144\), \(18\) of cardinality \(432\) and \(18\) of
cardinality \(1944\).
\end{proposicion}
\begin{proof}
The exhaustive evaluation of the algorithm in section~\ref{tpk:algoritmo}
over the indices \(0,\ldots,323\) gives the two preceding tables.
The first partition has \(18\) classes of cardinality \(18\), and the second
\(24\) classes of cardinality \(8\), one of cardinality \(24\) and one of
cardinality \(108\). The sums \(18\cdot18=324\) and
\(24\cdot8+24+108=324\) check the total mass of both partitions.
Proposition~\ref{tpk:recuperacion-firmas} makes the coupling
of their images injective; therefore, there are \(18\cdot26=468\) emissions.
By \eqref{tpk:factorizacion}, the preimage of a pair of signatures is
the product of their two preimages. One obtains \(18\cdot24=432\)
products of cardinality \(18\cdot8=144\), \(18\) of cardinality
\(18\cdot24=432\) and \(18\) of cardinality \(18\cdot108=1944\).
Finally,
\[
 432\cdot144+18\cdot432+18\cdot1944=104\,976.
\]
The enumerative part is an exact finite check of the recurrence;
the factorization explains why that count covers all pairs.
\end{proof}

The emission in example \ref{tpk:ejemplo501} has a fiber of cardinality
\(18\cdot24=432\). The example with indices \((0,0)\) has a fiber of
cardinality \(18\cdot8=144\). In both cases the emitted word recovers
the signatures, while identifying a particular initial condition
uses its index or its trajectory register.

\subsection{Oriented ternary reading and preservation of fibers}
\label{tpk:reduccion-ternaria}

The oriented reduction acts separately on each block:
\begin{equation}
 \Psi:\mathcal U_6\longrightarrow\mathbb F_3^6,\qquad
 \Psi(U_1,\ldots,U_6)=([-U_1]_3,\ldots,[-U_6]_3).
 \label{tpk:psi}
\end{equation}
Its sign fixes an orientation of the catalog. The balanced identification
\(0\mapsto0\), \(1\mapsto+1\), \(2\mapsto-1\) is applied afterward,
component by component.

Modulus nine organizes the APP evaluations and positions; residues
modulo ten form each digit; base one thousand brings together three decimal positions;
modulus three determines the ternary signature of each block. These four uses
have defined objects and functions. In particular, \(\Psi\) reads six
blocks and produces six trits; a positional conversion of the integer formed
by eighteen digits would be another operation.

For the two calculated conditions one obtains
\begin{align*}
 \Psi(501,614,498,169,272,272)&=(0,1,0,2,1,1),\\
 \Psi(252,109,252,903,850,850)&=(0,2,0,0,2,2).
\end{align*}
The decimal emitter is defined before this reduction. The function of
\(\Psi\) is to produce the ternary object on which orbital
symmetry and regional selectors act. That function explains its position in the
genealogy: it links an already calculated emission to a structural classification.

The image \(\mathcal W_6:=\Psi(\mathcal U_6)\) contains exactly
\(243\) words within the ambient set of \(3^6=729\) words.
Enumeration by the final algorithm gives the histogram
\begin{equation}
\begin{array}{c|rrrrrr}
 |\Psi^{-1}(w)|&1&2&3&4&5&6\\\hline
 \text{number of words }w&132&66&18&3&6&18.
\end{array}
\label{tpk:fibras-ternarias}
\end{equation}
The sums of both readings are
\[
132+66+18+3+6+18=243,
\qquad132+2\cdot66+3\cdot18+4\cdot3+5\cdot6+6\cdot18=468.
\]
The first counts ternary words; the second recovers the number of
emissions distributed among their fibers.

\begin{ejemplo}[Two emissions with the same ternary signature]
\label{tpk:colision}
The emissions
\[
 (498,501,614,252,252,109),\qquad
 (870,810,923,252,252,109)
\]
are distinct and both have image \((0,0,1,0,0,2)\) under \(\Psi\).
Each has \(144\) initial conditions. The ternary word preserves
a reading class; the emissions and their preimages distinguish the data
that converge in that class.
\end{ejemplo}

Preservation of fibers is expressed concretely by
\[
 \mathcal F_w=\{\mathbf U\in\mathcal U_6:\Psi(\mathbf U)=w\},
 \qquad
 \Omega_{\mathbf U}=T_{\mathrm{em}}^{-1}(\mathbf U).
\]
Therefore, the complete preimage of \(w\) in the initial domain is the
disjoint union of the sets \(\Omega_{\mathbf U}\), with
\(\mathbf U\in\mathcal F_w\). The enriched register can preserve the
emission, its signatures and the initial condition together with its ternary reading.
The exact inverse then refers to the recorded datum with its type;
the isolated six-trit word maintains the precise scope of
\eqref{tpk:psi}.

\subsection{Dihedral action and the forty-three orbits}
\label{tpk:orbitas}

On six coordinates the permutations are defined
\begin{align}
 r(u_1,u_2,u_3,u_4,u_5,u_6)
   &=(u_3,u_1,u_2,u_5,u_6,u_4),\\
 s(u_1,u_2,u_3,u_4,u_5,u_6)
   &=(u_4,u_5,u_6,u_1,u_2,u_3).
 \label{tpk:generadores-diedrales}
\end{align}
The first rotates the two groups of three coordinates in conjugate directions;
the second interchanges the two groups. Their composition satisfies
\(r^3=s^2=1\) and \(srs=r^{-1}\). They thus generate an action of the dihedral group
\(D_3\), with six elements.

\begin{proposicion}[Equivariance and stability of the catalog]
\label{tpk:equivarianza}
The reduction \(\Psi\) commutes with the action of \(D_3\). The sets
\(\mathcal U_6\) and \(\mathcal W_6\) are stable under its generators.
\end{proposicion}
\begin{proof}
The equality \(\Psi(g\mathbf U)=g\Psi(\mathbf U)\) is verified in each
coordinate because \(g\) permutes positions and \(\Psi\) applies the same
reduction to each position. Catalog stability constitutes another
check: when generating its \(468\) elements through
\eqref{tpk:factorizacion}, one verifies that \(r\mathbf U\) and
\(s\mathbf U\) belong to that set again. The algorithm in
section~\ref{tpk:algoritmo} performs these membership checks.
Equivariance then transmits stability to its ternary image.
\end{proof}

\begin{proposicion}[Complete orbital decomposition]
\label{tpk:43}
The set \(\mathcal W_6\) decomposes into \(43\) orbits: \(38\)
of cardinality six and \(5\) of cardinality three.
\end{proposicion}
\begin{proof}
For each \(w\in\mathcal W_6\) one forms
\[
 \mathcal O(w)=\{w,rw,r^2w,sw,srw,sr^2w\}.
\]
The lexicographically minimal word in that set is chosen as the
representative. Two words have the same representative exactly when
they belong to the same orbit. The algorithm generates the list of representatives
in table~\ref{tpk:representantes}; their cardinalities sum to
\(38\cdot6+5\cdot3=243\). Since every word in the catalog takes part in
this operation and the catalog is stable, the list covers the entire image.
\end{proof}

\begin{tablacompleta}[htbp]
\caption{Minimal representatives and cardinalities of the forty-three orbits.}
\label{tpk:representantes}
\small
\begin{tabular}{@{}lr@{\qquad}lr@{\qquad}lr@{}}
\toprule
representative&cardinality&representative&cardinality&representative&cardinality\\
\midrule
\texttt{001002}&6&\texttt{002112}&6&\texttt{012222}&6\\
\texttt{001011}&6&\texttt{002211}&6&\texttt{021022}&6\\
\texttt{001012}&6&\texttt{002220}&6&\texttt{021121}&6\\
\texttt{001021}&6&\texttt{011021}&6&\texttt{021202}&6\\
\texttt{001022}&6&\texttt{011112}&6&\texttt{022022}&3\\
\texttt{001100}&3&\texttt{011120}&6&\texttt{022112}&6\\
\texttt{001101}&6&\texttt{011201}&6&\texttt{022121}&6\\
\texttt{001110}&6&\texttt{011211}&6&\texttt{022202}&3\\
\texttt{001112}&6&\texttt{011220}&6&\texttt{022211}&6\\
\texttt{001202}&6&\texttt{011222}&6&\texttt{022222}&6\\
\texttt{001210}&6&\texttt{012112}&6&\texttt{112112}&3\\
\texttt{001211}&6&\texttt{012121}&6&\texttt{112211}&3\\
\texttt{001220}&6&\texttt{012202}&6&\texttt{112222}&6\\
\texttt{001222}&6&\texttt{012210}&6&&\\
\texttt{002012}&6&\texttt{012220}&6&&\\
\bottomrule
\end{tabular}
\end{tablacompleta}

Orbits inherit integer invariants from their fibers. For each word,
\(n(w)=|\mathcal F_w|\) counts emissions, while
\[
 M(w)=\sum_{\mathbf U\in\mathcal F_w}|\Omega_{\mathbf U}|
\]
counts initial conditions. The symbol census
\(c(w)=(n_0,n_1,n_2)\) records how many times each ternary class appears.
These data, together with the orientation and support of the paths,
provide the predicates for subsequent regional selection. The
classification thus preserves more information than the form of a representative
word.

\subsection{Self-contained enumeration and control algorithm}
\label{tpk:algoritmo}

The enumeration uses integers, finite lists and the operations already defined.
The following procedure specifies its execution; \(\sigma\) takes the values
\(+\) and \(\times\).

\begin{enumerate}
\item For each \(\nu=0,\ldots,323\), recover \((i,j;d)\) through
the inverse of \eqref{tpk:indice-semilla}. Initialize an empty list of signatures.
\item For each window \(m=1,\ldots,6\), initialize the accumulator \(A=0\).
Traverse \(t=9m-8,\ldots,9m\), calculating \(\epsilon_t\) by
\eqref{tpk:fase}.
\item If \(\sigma=+\) and \(\epsilon_t=+1\), add
\(\rho_9(i+j+2)\) to \(A\) and update
\((i,j)\leftarrow T_d(i,j)\). If \(\sigma=\times\) and
\(\epsilon_t=-1\), add
\(\ell_9(\rho_9((i+1)(j+1)))\) and perform the same transport.
\item Upon completing \(t=27\) or \(t=54\), replace \(d\) by its opposite.
Upon concluding the window, add \([A]_{10}\) to the signature list.
Preserve the position and direction for the next window.
\item Group the \(324\) indices by equality of their six-component
lists. These two groupings give \(f_+\), \(f_\times\) and all
their preimages.
\item For each pair of distinct signatures, apply \(G^{(6)}\).
Assign to the emission the product multiplicity of the two preimages.
Check signature recovery in each block and the sum of
multiplicities \(104\,976\).
\item Apply \(\Psi\) to each emission and group by equality of words.
Count the \(243\) values and the multiplicities of
\eqref{tpk:fibras-ternarias}.
\item Apply \(r,s\) to each emission and verify membership in the catalog.
For each ternary word, form its six dihedral images, remove
repetitions and record its lexicographic minimum. Count the \(43\)
orbits and the two sizes in table~\ref{tpk:representantes}.
\end{enumerate}

Execution terminates: the first stage evaluates \(648\) trajectories of
\(54\) transitions; the second combines \(468\) pairs of signatures; the
third groups a set of \(243\) words. An additional direct
check can execute the \(104\,976\) pairs and check the
factorization \eqref{tpk:factorizacion}; its scope coincides with the same
finite domain.

\subsection{Preservation of the register and function of the seeds}
\label{tpk:conclusion}

Let \(\mathsf{Sample}(q_t)\) be the sample preceding the transition and
\(\mathcal M_t\) the accumulated register. The memory update is
\begin{equation}
 \mathcal M_{t+1}=\mathcal M_t\mathbin{\Vert}\mathsf{Sample}(q_t),
 \label{tpk:archivo}
\end{equation}
where \(\Vert\) denotes ordered concatenation. Induction gives
\[
 \mathcal M_n=\mathcal M_0\mathbin{\Vert}
 (\mathsf{Sample}(q_0),\ldots,\mathsf{Sample}(q_{n-1})).
\]
The accumulators are reconstructed by summing the active contributions of
each segment of nine records. Their quotients, the cursor turns and
the orientations remain associated with the same transitions.
Execution determines that register through the recurrence; the archive
preserves the specific provenance of the readings.

At the end of this stage, the complete domain of initial
conditions, the decimal emission, its two recoverable signatures, the ternary image,
its fibers and its orbital classification have been constructed. The censuses \(104\,976\), \(468\),
\(243\), \(729\) and \(43\) now have a domain and a mechanism of obtainment.
Regional selections will receive these structures with their invariants,
and lifts will act on oriented words with their provenance
preserved. That continuity makes the finite catalog a constructive
antecedent of the coinductive development.
